
\chapter{Le problème isopérimétrique et la constante de Cheeger} \label{sec:cheeger_constant}

Cette annexe a pour objectif d'introduire la notion de \defini{constante de Cheeger} d'une variété compacte (une telle notion a été introduite pour les graphes dans la \Cref{chap:graphs}) et de fournir des éléments permettant d'aboutir à la conclusion que les minimiseurs de Cheeger existent pour toute variété compacte. Le contenu de cette annexe provient de \cite[Section 3]{benson_cheeger_2016}.

\begin{remarque}
Dans le \Cref{chap:demo_ineg_poincare}, nous utilisons la notation $\mesure$ pour désigner la mesure d'une partie de $\R^3$ (voir la \Cref{sec:triangulations}). Dans cette annexe, nous utilisons les notations $\vol_{n-1}$ (resp. $\vol_n$) pour désigner la mesure de dimension $n - 1$ (resp. $n$).
\end{remarque}

\begin{remarque}
    Les résultats suivants sont énoncés en dimension $n \geqslant 1$, mais seront ensuite appliqués au cas $n = 2$, qui est celui qui nous intéresse.
\end{remarque}

Soit $\surflisse$ une variété riemannienne de dimension $n$ telle que $\vol_n(\surflisse) < \infty$. La \defini{constante de Cheeger} de $\surflisse$ est définie par
\[
h(\surflisse) \defeq \inf_D \frac{\vol_{n - 1}(\partial D)}{\vol_n(D)},
\]
où $D \subset \surflisse$ est une sous-variété lisse de dimension $n$, à bord, telle que $0 < \vol_n(D) \leqslant \vol_n(\surflisse) / 2$.

Pour un sous-ensemble $A \subseteq \surflisse$ de mesure de Hausdorff $n$-dimensionnelle strictement positive et dont le bord $\partial A$ est mesurable pour la mesure de Hausdorff $(n-1)$-dimensionnelle, le \defini{rapport isopérimétrique} de $A$ est donné par
\begin{equation}
    h^\ast(A) \defeq \frac{\vol_{n-1}(\partial A)}{\vol_n(A)}.
\end{equation}
Le rapport isopérimétrique est bien défini lorsque $A$ est une sous-variété lisse à bord de $\surflisse$ ou, plus généralement, lorsque $A$ est un courant intégral. Nous définissons également la \defini{constante isopérimétrique} d'un sous-ensemble $A \subseteq \surflisse$ par
\begin{equation}
\overline{h}(A) \defeq \inf_{D \subseteq A} h^\ast(D),
\end{equation}
où $D$ parcourt tous les sous-ensembles de $A$ dont le bord est lisse.

Pour $t \in \intervalleoo{0}{\vol_n(\surflisse)}$, le problème isopérimétrique de volume $t$ consiste à rechercher un sous-ensemble $A$ de $\surflisse$ tel que $\vol_n(A) = t$ et
\[
\vol_{n-1}(\partial A) = \inf \ens[\big]{\vol_{n-1}(\partial B) \tq \vol_n(B) = t}.
\]
Un tel ensemble $A$ est appelé un \defini{minimiseur isopérimétrique}. Pour une présentation approfondie du problème isopérimétrique sur les variétés riemanniennes, voir \cite{ros_isoperimetric_2005}.

La constante de Cheeger est reliée au problème isopérimétrique en ce que nous cherchons à minimiser le rapport isopérimétrique parmi tous les problèmes isopérimétriques correspondant aux volumes $t \in \intervalleof{0}{\vol_n(\surflisse)/2}$. Nous appellerons \defini{candidat de Cheeger} tout ensemble $D$ de $\surflisse$ pour lequel $h^\ast(D)$ est bien défini et tel que $\vol_n(D) \leqslant \vol_n(\surflisse)/2$.

Comme indiqué dans \cite{ros_isoperimetric_2005}, on obtient le théorème fondamental d'existence et de régularité pour les minimiseurs isopérimétriques des variétés riemanniennes.

\begin{theo}[Existence et régularité] \label{theo:fundamental_existence_and_regularity}
    Si $\surflisse$ est une variété lisse, compacte, éventuellement à bord, et si $t \in \intervalleoo{0}{\vol_n(\surflisse)}$, alors il existe une région compacte $A \subset \surflisse$ telle que $\partial A$ minimise le volume $(n-1)$-dimensionnel parmi toutes les régions de volume $n$-dimensionnel égal à $t$. De plus, pour de tels minimiseurs, $\partial A$ est une hypersurface lisse plongée, de courbure moyenne constante, à l'exception d'un ensemble singulier de dimension de Hausdorff au plus égale à $n - 8$. Enfin, si $\partial \surflisse \cap \partial A \not= \varnothing$, alors $\partial \surflisse$ et $\partial A$ se rencontrent orthogonalement.
\end{theo}

Comme le suggèrent le \Cref{theo:fundamental_existence_and_regularity} et les techniques d'approximation lisse, la recherche de la constante de Cheeger peut se faire de manière équivalente en minimisant le rapport isopérimétrique sur les sous-variétés lisses à bord ou sur la classe plus générale des courants rectifiables. De même que nous avons défini les minimiseurs isopérimétriques, nous pouvons définir les minimiseurs de la constante de Cheeger, que nous appellerons simplement minimiseurs de Cheeger. Plus précisément, un \defini{minimiseur de Cheeger} est un sous-ensemble $A$ de $\surflisse$ tel que $\vol_n(A) \leqslant \vol_n(\surflisse)/2$ et $h^\ast(A) = h(\surflisse)$. Lorsque $\surflisse$ est compacte, Buser a montré qu'il existe un ou plusieurs volumes $t \in \intervalleof{0}{\vol_n(\surflisse)/2}$ pour lesquels les minimiseurs de Cheeger coïncident avec les solutions du problème isopérimétrique de volume $n$-dimensionnel $t$ dans $\surflisse$ \cite[Remark 3.3, Lemma 3.4]{buser_note_1982}. L'existence et la régularité d'un minimiseur de Cheeger découlent donc du \Cref{theo:fundamental_existence_and_regularity}. Plus précisément, Buser a démontré le résultat suivant.

\begin{lemme}[Buser, {\cite[Lemme 3.4]{buser_note_1982}}] \label{lemme:buser}
Soit $\surflisse$ une variété riemannienne compacte de dimension $n$, éventuellement à bord. Alors il existe $t^\ast \in \intervalleof{0}{\vol_n(\surflisse)/2}$ et une suite de sous-variétés $A_k \subset \surflisse$ à bord lisse telles que $\vol_n(A_k) = t^\ast$ pour tout $k$ et
\[
h(\surflisse) = \lim_{k \to \infty} \frac{\vol_{n-1}(\partial A_k)}{\vol_n(A_k)}.
\]
\end{lemme}
En combinant le \Cref{theo:fundamental_existence_and_regularity} et le \Cref{lemme:buser}, nous concluons immédiatement que les minimiseurs de Cheeger existent pour toute variété compacte.